\chapter{Matrix Multiplication}

Matrix multiplication is \emph{the} central operation of linear algebra: it
is how linear maps compose, how systems of equations are written compactly
as $A x = b$, and how almost every formula in this sheet (from
$\hat{\beta} = (X^{T} X)^{-1} X^{T} y$ to $A^{n} = P D^{n} P^{-1}$) is
built. Its conventions take some getting used to: dimensions must mesh, and
order matters.

\section{Core Rules}

\begin{rulebox}{Definition}
For $A \in \R^{m \times n}$ and $B \in \R^{n \times p}$, the product
$C = A B \in \R^{m \times p}$ has entries
\[
c_{ij} \;=\; \sum_{k=1}^{n} a_{ik} \, b_{kj}, \qquad 1 \le i \le m, \; 1 \le j \le p.
\]
That is, the $(i, j)$-entry of $A B$ is the dot product of row $i$ of $A$
with column $j$ of $B$.
\end{rulebox}

\begin{rulebox}{Dimension compatibility}
The product $A B$ is defined only when the number of columns of $A$ equals
the number of rows of $B$: $(m \times n)(n \times p) = (m \times p)$. The
``inner'' dimensions must match; the ``outer'' dimensions give the result.
\end{rulebox}

\begin{rulebox}{Non-commutativity}
In general $A B \neq B A$, even when both products are defined.
\end{rulebox}

\begin{rulebox}{Algebraic rules}
Where the products are defined:
\begin{itemize}[leftmargin=1.6em,itemsep=0.3ex]
    \item \textbf{Associativity:} $(A B) C = A (B C)$.
    \item \textbf{Distributivity:} $A (B + C) = A B + A C$ and $(A + B) C = A C + B C$.
    \item \textbf{Scalar compatibility:} $(\lambda A) B = \lambda (A B) = A (\lambda B)$.
    \item \textbf{Identity:} $I_{m} A = A = A I_{n}$ for $A \in \R^{m \times n}$.
    \item \textbf{Transpose of a product:} $(A B)^{T} = B^{T} A^{T}$.
\end{itemize}
\end{rulebox}

\begin{rulebox}{Matrix--vector products}
For $A \in \R^{m \times n}$ and $x \in \R^{n}$, $A x \in \R^{m}$ is a linear
combination of the columns of $A$ with coefficients $x_{1}, \dots, x_{n}$.
\end{rulebox}

\begin{rulebox}{Powers}
For a square matrix $A \in \R^{n \times n}$, $A^{0} = I$, $A^{1} = A$, and
$A^{k} = A \cdot A^{k - 1}$ for $k \ge 1$.
\end{rulebox}

\section{Worked Example}

\begin{worked}{With $A=\begin{pmatrix}1&2&3\\0&-1&4\end{pmatrix}$ and $B=\begin{pmatrix}2&1\\0&-1\\1&3\end{pmatrix}$, compute $AB$ and $BA$; then show that $A^{2}=0$ does not force $A=0$}
\wstep{Step 1. Shapes first: they tell you the answer's size, or that there is no answer.}
$A$ is $2\times3$, $B$ is $3\times2$. Write the shapes side by side:
\[
  \underbrace{(2\times\mathbf 3)}_{A}\underbrace{(\mathbf 3\times2)}_{B}
  \;\Longrightarrow\;2\times2,
  \qquad
  \underbrace{(3\times\mathbf 2)}_{B}\underbrace{(\mathbf 2\times3)}_{A}
  \;\Longrightarrow\;3\times3 .
\]
The inner numbers must agree and then cancel; the outer numbers survive as
the result's shape. So both products exist here, and they are not even the
same size --- non-commutativity in its starkest form.

\wstep{Step 2. Compute $AB$: entry $(i,j)$ is row $i$ of $A$ dotted with column $j$ of $B$.}
\[
  (AB)_{11}=1\cdot2+2\cdot0+3\cdot1=5,\qquad
  (AB)_{12}=1\cdot1+2\cdot(-1)+3\cdot3=8,
\]
\[
  (AB)_{21}=0\cdot2+(-1)\cdot0+4\cdot1=4,\qquad
  (AB)_{22}=0\cdot1+(-1)(-1)+4\cdot3=13 .
\]
\[
  AB=\begin{pmatrix}5&8\\4&13\end{pmatrix}
\]
\wtrap{\emph{Row} of the left factor with \emph{column} of the right. Taking
row-with-row is the standard beginner's error and it silently produces a
plausible-looking wrong matrix. Say ``row $i$, column $j$'' out loud as you
compute each entry.}

\wstep{Step 3. Compute $BA$, which is a $3\times3$ matrix.}
Row by row of $B$ against the three columns of $A$, which are
$(1,0)$, $(2,-1)$, $(3,4)$:
\[
  \text{row }(2,1):\ (2,\;3,\;10);\qquad
  \text{row }(0,-1):\ (0,\;1,\;-4);\qquad
  \text{row }(1,3):\ (1,\;-1,\;15).
\]
\[
  BA=\begin{pmatrix}2&3&10\\0&1&-4\\1&-1&15\end{pmatrix}
\]

\wstep{Step 4. Read the moral.}
$AB$ is $2\times2$ and $BA$ is $3\times3$. \textbf{Matrix multiplication is
not commutative} --- here they are not even comparable. Even for square
matrices of the same size, $AB\neq BA$ is the rule and equality the
exception.
\wtrap{Because order matters, every algebraic manipulation must preserve it.
You may not ``cancel'' $A$ from $AB=AC$ unless $A$ is invertible, and
$(A+B)^{2}=A^{2}+AB+BA+B^{2}$, which equals $A^{2}+2AB+B^{2}$ only when $A$
and $B$ commute.}

\wstep{Step 5. A check that costs one line: the trace.}
For any $A,B$ of compatible shapes, $\operatorname{tr}(AB)=\operatorname{tr}(BA)$
even though the matrices differ. Here $\operatorname{tr}(AB)=5+13=18$ and
$\operatorname{tr}(BA)=2+1+15=18$. \checkmark\ Two independent computations
agreeing on one number catches most arithmetic slips.

\wstep{Step 6. Now the second surprise: a non-zero matrix whose square is zero.}
Take
\[
  N=\begin{pmatrix}0&1\\0&0\end{pmatrix},
  \qquad
  N^{2}=\begin{pmatrix}0&1\\0&0\end{pmatrix}\begin{pmatrix}0&1\\0&0\end{pmatrix}
  =\begin{pmatrix}0&0\\0&0\end{pmatrix} .
\]
So $N^{2}=0$ with $N\neq0$. There is \textbf{no zero-product property} for
matrices: $MN=0$ tells you nothing about $M$ or $N$ individually.
\wtrap{This is what makes matrix algebra genuinely different from the
algebra of numbers, and it has consequences everywhere: $A^{2}=A$ does not
imply $A\in\{0,I\}$ (projection matrices satisfy it), and from $AB=0$ you
may not conclude that one factor vanishes. When you are tempted to divide,
stop and ask whether an inverse exists.}

\wstep{Step 7. The reading that makes products intuitive.}
$A\mathbf x$ is a \textbf{linear combination of the columns of $A$}, with
the entries of $\mathbf x$ as weights. For instance
\[
  A\begin{pmatrix}2\\0\\1\end{pmatrix}
  =2\begin{pmatrix}1\\0\end{pmatrix}
  +0\begin{pmatrix}2\\-1\end{pmatrix}
  +1\begin{pmatrix}3\\4\end{pmatrix}
  =\begin{pmatrix}5\\4\end{pmatrix},
\]
which is exactly the first column of $AB$ --- because the first column of
$B$ is $(2,0,1)^{\!\top}$. Every column of a product is the left matrix
acting on the corresponding column of the right one, and that single
observation explains rank, column space, and solvability in Chapter~7.
\end{worked}

\section{Practice Problems}

Throughout, let
\[
A = \begin{pmatrix} 1 & 2 \\ 3 & 4 \end{pmatrix}, \quad
B = \begin{pmatrix} 0 & -1 \\ 2 & 5 \end{pmatrix}, \quad
P = \begin{pmatrix} 1 & 0 & 2 \\ -1 & 3 & 1 \end{pmatrix}, \quad
Q = \begin{pmatrix} 1 & 2 \\ 0 & -1 \\ 3 & 1 \end{pmatrix}, \quad x = \begin{pmatrix} 1 \\ 2 \end{pmatrix}.
\]

\begin{problems}
    \item Compute $A B$.
    \item Compute $B A$. Compare with $A B$.
    \item Compute $A x$.
    \item Compute $B x$.
    \item Compute $A^{2}$.
    \item Compute $A^{3}$.
    \item Compute $P Q$.
    \item Compute $Q P$.
    \item Which of $P P$, $Q Q$, $P Q$, $Q P$ are defined? State the dimensions of each defined product.
    \item Compute $I_{2} \, A$ and check it equals $A$.
    \item Compute $(A B)^{T}$ and verify it equals $B^{T} A^{T}$.
    \item Let $u = (1, -1)^{T}$, $v = (2, 3)^{T}$. Compute $u^{T} v$ and $u v^{T}$.
    \item For $A$ as above, compute $A^{T} A$.
    \item Is $A^{T} A$ symmetric? Justify (general statement).
    \item Find a $2 \times 2$ matrix $M \neq 0$ such that $M^{2} = 0$.
    \item Find $2 \times 2$ matrices $A, B$ with $A B = 0$ but neither $A = 0$ nor $B = 0$.
    \item Compute $\bigl(\begin{smallmatrix} 1 & 1 \\ 0 & 1 \end{smallmatrix}\bigr)^{n}$ for general $n \ge 1$ (conjecture and verify for $n = 2, 3$).
    \item For $D = \mathrm{diag}(2, 3)$, compute $D^{4}$.
    \item For $A$ above, compute $A (B + I_{2})$ in two ways: directly, and as $A B + A$.
    \item Let $e_{1} = (1, 0)^{T}$, $e_{2} = (0, 1)^{T}$. Compute $A e_{1}$ and $A e_{2}$, and observe they are the columns of $A$.
    \item Let $J = \bigl(\begin{smallmatrix} 0 & -1 \\ 1 & 0 \end{smallmatrix}\bigr)$. Compute $J^{2}$.
    \item For $J$ as above, compute $J^{4}$.
    \item Let $R = \bigl(\begin{smallmatrix} 1 & 1 \\ 1 & 1 \end{smallmatrix}\bigr)$. Compute $R^{2}$ and conjecture $R^{n}$.
    \item Compute $(A + B)(A - B)$ for $A, B$ above. Is it equal to $A^{2} - B^{2}$?
    \item For an upper-triangular $2 \times 2$ matrix $U = \bigl(\begin{smallmatrix} a & b \\ 0 & c \end{smallmatrix}\bigr)$, compute $U^{2}$ and state when the result is again upper-triangular.
\end{problems}

\section{Solutions}

\begin{solutions}
    \item $A B = \begin{pmatrix} 1 \cdot 0 + 2 \cdot 2 & 1 \cdot (-1) + 2 \cdot 5 \\ 3 \cdot 0 + 4 \cdot 2 & 3 \cdot (-1) + 4 \cdot 5 \end{pmatrix} = \begin{pmatrix} 4 & 9 \\ 8 & 17 \end{pmatrix}$.

    \item $B A = \begin{pmatrix} 0 \cdot 1 + (-1) \cdot 3 & 0 \cdot 2 + (-1) \cdot 4 \\ 2 \cdot 1 + 5 \cdot 3 & 2 \cdot 2 + 5 \cdot 4 \end{pmatrix} = \begin{pmatrix} -3 & -4 \\ 17 & 24 \end{pmatrix} \neq A B$.

    \item $A x = \begin{pmatrix} 1 \cdot 1 + 2 \cdot 2 \\ 3 \cdot 1 + 4 \cdot 2 \end{pmatrix} = \begin{pmatrix} 5 \\ 11 \end{pmatrix}$.

    \item $B x = \begin{pmatrix} 0 \cdot 1 + (-1) \cdot 2 \\ 2 \cdot 1 + 5 \cdot 2 \end{pmatrix} = \begin{pmatrix} -2 \\ 12 \end{pmatrix}$.

    \item $A^{2} = A \cdot A = \begin{pmatrix} 1 + 6 & 2 + 8 \\ 3 + 12 & 6 + 16 \end{pmatrix} = \begin{pmatrix} 7 & 10 \\ 15 & 22 \end{pmatrix}$.

    \item $A^{3} = A \cdot A^{2} = \begin{pmatrix} 1 & 2 \\ 3 & 4 \end{pmatrix}\begin{pmatrix} 7 & 10 \\ 15 & 22 \end{pmatrix} = \begin{pmatrix} 7 + 30 & 10 + 44 \\ 21 + 60 & 30 + 88 \end{pmatrix} = \begin{pmatrix} 37 & 54 \\ 81 & 118 \end{pmatrix}$.

    \item $P Q$: $P \in \R^{2 \times 3}$, $Q \in \R^{3 \times 2}$, so $P Q \in \R^{2 \times 2}$.
    \[
    P Q = \begin{pmatrix} 1 \cdot 1 + 0 \cdot 0 + 2 \cdot 3 & 1 \cdot 2 + 0 \cdot (-1) + 2 \cdot 1 \\ -1 \cdot 1 + 3 \cdot 0 + 1 \cdot 3 & -1 \cdot 2 + 3 \cdot (-1) + 1 \cdot 1 \end{pmatrix} = \begin{pmatrix} 7 & 4 \\ 2 & -4 \end{pmatrix}.
    \]

    \item $Q P$: $Q \in \R^{3 \times 2}$, $P \in \R^{2 \times 3}$, so $Q P \in \R^{3 \times 3}$.
    \[
    Q P = \begin{pmatrix} 1 \cdot 1 + 2 \cdot (-1) & 1 \cdot 0 + 2 \cdot 3 & 1 \cdot 2 + 2 \cdot 1 \\ 0 \cdot 1 + (-1)(-1) & 0 \cdot 0 + (-1) \cdot 3 & 0 \cdot 2 + (-1) \cdot 1 \\ 3 \cdot 1 + 1 \cdot (-1) & 3 \cdot 0 + 1 \cdot 3 & 3 \cdot 2 + 1 \cdot 1 \end{pmatrix} = \begin{pmatrix} -1 & 6 & 4 \\ 1 & -3 & -1 \\ 2 & 3 & 7 \end{pmatrix}.
    \]

    \item $P P$ undefined ($2 \times 3$ times $2 \times 3$). $Q Q$ undefined ($3 \times 2$ times $3 \times 2$). $P Q \in \R^{2 \times 2}$. $Q P \in \R^{3 \times 3}$.

    \item $I_{2} A = A$ by direct multiplication (each row of $I_{2}$ picks out the corresponding row of $A$).

    \item $(A B)^{T} = \bigl(\begin{smallmatrix} 4 & 9 \\ 8 & 17 \end{smallmatrix}\bigr)^{T} = \bigl(\begin{smallmatrix} 4 & 8 \\ 9 & 17 \end{smallmatrix}\bigr)$. And $B^{T} A^{T} = \bigl(\begin{smallmatrix} 0 & 2 \\ -1 & 5 \end{smallmatrix}\bigr)\bigl(\begin{smallmatrix} 1 & 3 \\ 2 & 4 \end{smallmatrix}\bigr) = \bigl(\begin{smallmatrix} 4 & 8 \\ 9 & 17 \end{smallmatrix}\bigr)$. Equal.

    \item $u^{T} v = 1 \cdot 2 + (-1) \cdot 3 = -1$ (a scalar). $u v^{T} = \bigl(\begin{smallmatrix} 1 \\ -1 \end{smallmatrix}\bigr)\begin{pmatrix} 2 & 3 \end{pmatrix} = \bigl(\begin{smallmatrix} 2 & 3 \\ -2 & -3 \end{smallmatrix}\bigr)$ (a rank-one $2 \times 2$ matrix).

    \item $A^{T} A = \bigl(\begin{smallmatrix} 1 & 3 \\ 2 & 4 \end{smallmatrix}\bigr)\bigl(\begin{smallmatrix} 1 & 2 \\ 3 & 4 \end{smallmatrix}\bigr) = \bigl(\begin{smallmatrix} 1 + 9 & 2 + 12 \\ 2 + 12 & 4 + 16 \end{smallmatrix}\bigr) = \bigl(\begin{smallmatrix} 10 & 14 \\ 14 & 20 \end{smallmatrix}\bigr)$.

    \item Yes: $(A^{T} A)^{T} = A^{T} (A^{T})^{T} = A^{T} A$. So $A^{T} A$ is always symmetric.

    \item E.g.\ $M = \bigl(\begin{smallmatrix} 0 & 1 \\ 0 & 0 \end{smallmatrix}\bigr)$. Then $M^{2} = \bigl(\begin{smallmatrix} 0 & 0 \\ 0 & 0 \end{smallmatrix}\bigr)$.

    \item E.g.\ $A = \bigl(\begin{smallmatrix} 1 & 0 \\ 0 & 0 \end{smallmatrix}\bigr)$, $B = \bigl(\begin{smallmatrix} 0 & 0 \\ 0 & 1 \end{smallmatrix}\bigr)$. Then $A B = 0$.

    \item Let $U = \bigl(\begin{smallmatrix} 1 & 1 \\ 0 & 1 \end{smallmatrix}\bigr)$. Then $U^{2} = \bigl(\begin{smallmatrix} 1 & 2 \\ 0 & 1 \end{smallmatrix}\bigr)$, $U^{3} = \bigl(\begin{smallmatrix} 1 & 3 \\ 0 & 1 \end{smallmatrix}\bigr)$. Conjecture: $U^{n} = \bigl(\begin{smallmatrix} 1 & n \\ 0 & 1 \end{smallmatrix}\bigr)$ (proof by induction: $U^{n+1} = U^{n} U$).

    \item $D^{n} = \mathrm{diag}(2^{n}, 3^{n})$, so $D^{4} = \mathrm{diag}(16, 81) = \bigl(\begin{smallmatrix} 16 & 0 \\ 0 & 81 \end{smallmatrix}\bigr)$.

    \item $B + I_{2} = \bigl(\begin{smallmatrix} 1 & -1 \\ 2 & 6 \end{smallmatrix}\bigr)$. Direct: $A(B + I_{2}) = \bigl(\begin{smallmatrix} 1 & 2 \\ 3 & 4 \end{smallmatrix}\bigr)\bigl(\begin{smallmatrix} 1 & -1 \\ 2 & 6 \end{smallmatrix}\bigr) = \bigl(\begin{smallmatrix} 5 & 11 \\ 11 & 21 \end{smallmatrix}\bigr)$. Check: $A B + A = \bigl(\begin{smallmatrix} 4 & 9 \\ 8 & 17 \end{smallmatrix}\bigr) + \bigl(\begin{smallmatrix} 1 & 2 \\ 3 & 4 \end{smallmatrix}\bigr) = \bigl(\begin{smallmatrix} 5 & 11 \\ 11 & 21 \end{smallmatrix}\bigr)$. Match.

    \item $A e_{1} = \bigl(\begin{smallmatrix} 1 \\ 3 \end{smallmatrix}\bigr)$, $A e_{2} = \bigl(\begin{smallmatrix} 2 \\ 4 \end{smallmatrix}\bigr)$, which are the two columns of $A$.

    \item $J^{2} = \bigl(\begin{smallmatrix} 0 & -1 \\ 1 & 0 \end{smallmatrix}\bigr)\bigl(\begin{smallmatrix} 0 & -1 \\ 1 & 0 \end{smallmatrix}\bigr) = \bigl(\begin{smallmatrix} -1 & 0 \\ 0 & -1 \end{smallmatrix}\bigr) = -I_{2}$. ($J$ models multiplication by $i$.)

    \item $J^{4} = (J^{2})^{2} = (-I)^{2} = I$.

    \item $R^{2} = \bigl(\begin{smallmatrix} 1 & 1 \\ 1 & 1 \end{smallmatrix}\bigr)\bigl(\begin{smallmatrix} 1 & 1 \\ 1 & 1 \end{smallmatrix}\bigr) = \bigl(\begin{smallmatrix} 2 & 2 \\ 2 & 2 \end{smallmatrix}\bigr) = 2 R$. By induction $R^{n} = 2^{n - 1} R$ for $n \ge 1$.

    \item $A + B = \bigl(\begin{smallmatrix} 1 & 1 \\ 5 & 9 \end{smallmatrix}\bigr)$, $A - B = \bigl(\begin{smallmatrix} 1 & 3 \\ 1 & -1 \end{smallmatrix}\bigr)$, product
    \[
    (A + B)(A - B) = \bigl(\begin{smallmatrix} 1 + 1 & 3 - 1 \\ 5 + 9 & 15 - 9 \end{smallmatrix}\bigr) = \bigl(\begin{smallmatrix} 2 & 2 \\ 14 & 6 \end{smallmatrix}\bigr).
    \]
    Now $A^{2} = \bigl(\begin{smallmatrix} 7 & 10 \\ 15 & 22 \end{smallmatrix}\bigr)$ and $B^{2} = \bigl(\begin{smallmatrix} 0 & -1 \\ 2 & 5 \end{smallmatrix}\bigr)\bigl(\begin{smallmatrix} 0 & -1 \\ 2 & 5 \end{smallmatrix}\bigr) = \bigl(\begin{smallmatrix} -2 & -5 \\ 10 & 23 \end{smallmatrix}\bigr)$, so $A^{2} - B^{2} = \bigl(\begin{smallmatrix} 9 & 15 \\ 5 & -1 \end{smallmatrix}\bigr) \neq (A+B)(A-B)$. The reason: $(A+B)(A-B) = A^{2} - A B + B A - B^{2}$, and the cross terms $-A B + B A$ vanish only when $A B = B A$.

    \item $U^{2} = \bigl(\begin{smallmatrix} a & b \\ 0 & c \end{smallmatrix}\bigr)\bigl(\begin{smallmatrix} a & b \\ 0 & c \end{smallmatrix}\bigr) = \bigl(\begin{smallmatrix} a^{2} & a b + b c \\ 0 & c^{2} \end{smallmatrix}\bigr)$. The $(2, 1)$-entry is $0$, so $U^{2}$ is always upper-triangular. (Same is true for any product of upper-triangular matrices.)
\end{solutions}

\section{Additional Problems}

Throughout, let
\[
E = \begin{pmatrix} 2 & -1 \\ 1 & 3 \end{pmatrix}, \quad
F = \begin{pmatrix} 1 & 0 \\ -2 & 4 \end{pmatrix}, \quad
M = \begin{pmatrix} 1 & 2 & -1 \\ 0 & 1 & 3 \end{pmatrix}, \quad
N = \begin{pmatrix} 2 & 1 \\ 1 & 0 \\ -1 & 2 \end{pmatrix}, \quad y = \begin{pmatrix} 2 \\ -1 \end{pmatrix}.
\]

\begin{problems}
    \item Compute $E F$.
    \item Compute $F E$. Compare with $E F$.
    \item Compute $E y$.
    \item Compute $F y$.
    \item Compute $E^{2}$.
    \item Compute $F^{2}$.
    \item Compute $M N$.
    \item Compute $N M$.
    \item Which of $M M$, $N N$, $M N$, $N M$ are defined? Give the dimensions.
    \item Compute $(E F)^{T}$ and verify it equals $F^{T} E^{T}$.
    \item Let $u = (2, 1)^{T}$, $v = (1, 3)^{T}$. Compute $u^{T} v$ and $v u^{T}$.
    \item For $E$ above, compute $E E^{T}$.
    \item Is $E E^{T}$ symmetric? Justify.
    \item Find a nonzero $2 \times 2$ matrix $K$ with $K^{2} = K$ (idempotent, $\neq I$).
    \item Compute $\bigl(\begin{smallmatrix} 2 & 0 \\ 0 & -1 \end{smallmatrix}\bigr)^{n}$ for general $n \ge 1$.
    \item For $D = \mathrm{diag}(3, -2, 5)$, compute $D^{3}$.
    \item For $E$ above, compute $E(F + I_{2})$ in two ways: directly and as $E F + E$.
    \item Let $e_{1} = (1, 0)^{T}$, $e_{2} = (0, 1)^{T}$. Compute $E^{T} e_{1}$ and $E^{T} e_{2}$ and identify the result.
    \item Let $S = \bigl(\begin{smallmatrix} 0 & 1 \\ 1 & 0 \end{smallmatrix}\bigr)$. Compute $S^{2}$.
    \item Let $T = \bigl(\begin{smallmatrix} 1 & 2 \\ 3 & 4 \end{smallmatrix}\bigr)$ and find $S T$ and $T S$ (with $S$ as above). What do these represent?
    \item Let $R = \bigl(\begin{smallmatrix} 2 & 2 \\ 2 & 2 \end{smallmatrix}\bigr)$. Compute $R^{2}$ and conjecture $R^{n}$ for $n \ge 1$.
    \item Compute $(E - F)(E + F)$ directly and compare with $E^{2} - F^{2}$.
    \item For lower-triangular $L = \bigl(\begin{smallmatrix} a & 0 \\ b & c \end{smallmatrix}\bigr)$, compute $L^{2}$ and confirm it is lower-triangular.
    \item Find $2 \times 2$ matrices $A, B$ with $A B = I$ but $B A \neq I$. (Try $2 \times 2$ -- impossible! Prove it.) -- show that for $2 \times 2$ matrices, $A B = I \Rightarrow B A = I$.
    \item Compute $N^{T} N$ for $N$ above.
\end{problems}

\section{Additional Solutions}

\begin{solutions}
    \item $E F = \bigl(\begin{smallmatrix} 2 \cdot 1 + (-1)(-2) & 2 \cdot 0 + (-1) \cdot 4 \\ 1 \cdot 1 + 3 \cdot (-2) & 1 \cdot 0 + 3 \cdot 4 \end{smallmatrix}\bigr) = \bigl(\begin{smallmatrix} 4 & -4 \\ -5 & 12 \end{smallmatrix}\bigr)$.

    \item $F E = \bigl(\begin{smallmatrix} 1 \cdot 2 + 0 \cdot 1 & 1 \cdot (-1) + 0 \cdot 3 \\ -2 \cdot 2 + 4 \cdot 1 & -2 \cdot (-1) + 4 \cdot 3 \end{smallmatrix}\bigr) = \bigl(\begin{smallmatrix} 2 & -1 \\ 0 & 14 \end{smallmatrix}\bigr) \neq E F$.

    \item $E y = \bigl(\begin{smallmatrix} 2 \cdot 2 + (-1)(-1) \\ 1 \cdot 2 + 3 \cdot (-1) \end{smallmatrix}\bigr) = \bigl(\begin{smallmatrix} 5 \\ -1 \end{smallmatrix}\bigr)$.

    \item $F y = \bigl(\begin{smallmatrix} 1 \cdot 2 + 0 \cdot (-1) \\ -2 \cdot 2 + 4 \cdot (-1) \end{smallmatrix}\bigr) = \bigl(\begin{smallmatrix} 2 \\ -8 \end{smallmatrix}\bigr)$.

    \item $E^{2} = \bigl(\begin{smallmatrix} 2 & -1 \\ 1 & 3 \end{smallmatrix}\bigr)\bigl(\begin{smallmatrix} 2 & -1 \\ 1 & 3 \end{smallmatrix}\bigr) = \bigl(\begin{smallmatrix} 4 - 1 & -2 - 3 \\ 2 + 3 & -1 + 9 \end{smallmatrix}\bigr) = \bigl(\begin{smallmatrix} 3 & -5 \\ 5 & 8 \end{smallmatrix}\bigr)$.

    \item $F^{2} = \bigl(\begin{smallmatrix} 1 & 0 \\ -2 & 4 \end{smallmatrix}\bigr)\bigl(\begin{smallmatrix} 1 & 0 \\ -2 & 4 \end{smallmatrix}\bigr) = \bigl(\begin{smallmatrix} 1 & 0 \\ -2 - 8 & 16 \end{smallmatrix}\bigr) = \bigl(\begin{smallmatrix} 1 & 0 \\ -10 & 16 \end{smallmatrix}\bigr)$.

    \item $M N$: $(2 \times 3)(3 \times 2) = 2 \times 2$.
    \[
    M N = \begin{pmatrix} 1 \cdot 2 + 2 \cdot 1 + (-1)(-1) & 1 \cdot 1 + 2 \cdot 0 + (-1) \cdot 2 \\ 0 \cdot 2 + 1 \cdot 1 + 3 \cdot (-1) & 0 \cdot 1 + 1 \cdot 0 + 3 \cdot 2 \end{pmatrix} = \begin{pmatrix} 5 & -1 \\ -2 & 6 \end{pmatrix}.
    \]

    \item $N M$: $(3 \times 2)(2 \times 3) = 3 \times 3$.
    \[
    N M = \begin{pmatrix} 2 \cdot 1 + 1 \cdot 0 & 2 \cdot 2 + 1 \cdot 1 & 2 \cdot (-1) + 1 \cdot 3 \\ 1 \cdot 1 + 0 \cdot 0 & 1 \cdot 2 + 0 \cdot 1 & 1 \cdot (-1) + 0 \cdot 3 \\ -1 \cdot 1 + 2 \cdot 0 & -1 \cdot 2 + 2 \cdot 1 & -1 \cdot (-1) + 2 \cdot 3 \end{pmatrix} = \begin{pmatrix} 2 & 5 & 1 \\ 1 & 2 & -1 \\ -1 & 0 & 7 \end{pmatrix}.
    \]

    \item $M M$: $2 \times 3$ times $2 \times 3$ undefined. $N N$: $3 \times 2$ times $3 \times 2$ undefined. $M N \in \R^{2 \times 2}$. $N M \in \R^{3 \times 3}$.

    \item $(E F)^{T} = \bigl(\begin{smallmatrix} 4 & -4 \\ -5 & 12 \end{smallmatrix}\bigr)^{T} = \bigl(\begin{smallmatrix} 4 & -5 \\ -4 & 12 \end{smallmatrix}\bigr)$. $F^{T} E^{T} = \bigl(\begin{smallmatrix} 1 & -2 \\ 0 & 4 \end{smallmatrix}\bigr)\bigl(\begin{smallmatrix} 2 & 1 \\ -1 & 3 \end{smallmatrix}\bigr) = \bigl(\begin{smallmatrix} 2 + 2 & 1 - 6 \\ 0 - 4 & 0 + 12 \end{smallmatrix}\bigr) = \bigl(\begin{smallmatrix} 4 & -5 \\ -4 & 12 \end{smallmatrix}\bigr)$. Equal.

    \item $u^{T} v = 2 \cdot 1 + 1 \cdot 3 = 5$. $v u^{T} = \bigl(\begin{smallmatrix} 1 \\ 3 \end{smallmatrix}\bigr)\begin{pmatrix} 2 & 1 \end{pmatrix} = \bigl(\begin{smallmatrix} 2 & 1 \\ 6 & 3 \end{smallmatrix}\bigr)$.

    \item $E E^{T} = \bigl(\begin{smallmatrix} 2 & -1 \\ 1 & 3 \end{smallmatrix}\bigr)\bigl(\begin{smallmatrix} 2 & 1 \\ -1 & 3 \end{smallmatrix}\bigr) = \bigl(\begin{smallmatrix} 4 + 1 & 2 - 3 \\ 2 - 3 & 1 + 9 \end{smallmatrix}\bigr) = \bigl(\begin{smallmatrix} 5 & -1 \\ -1 & 10 \end{smallmatrix}\bigr)$.

    \item Yes: $(E E^{T})^{T} = (E^{T})^{T} E^{T} = E E^{T}$. So $E E^{T}$ is always symmetric.

    \item E.g.\ $K = \bigl(\begin{smallmatrix} 1 & 0 \\ 0 & 0 \end{smallmatrix}\bigr)$: $K^{2} = K$. Or $K = \bigl(\begin{smallmatrix} 1 & 1 \\ 0 & 0 \end{smallmatrix}\bigr)$: $K^{2} = \bigl(\begin{smallmatrix} 1 & 1 \\ 0 & 0 \end{smallmatrix}\bigr) = K$.

    \item $D^{n} = \mathrm{diag}(2^{n}, (-1)^{n})$.

    \item $D^{3} = \mathrm{diag}(27, -8, 125)$.

    \item $F + I_{2} = \bigl(\begin{smallmatrix} 2 & 0 \\ -2 & 5 \end{smallmatrix}\bigr)$. Direct: $E(F + I_{2}) = \bigl(\begin{smallmatrix} 2 & -1 \\ 1 & 3 \end{smallmatrix}\bigr)\bigl(\begin{smallmatrix} 2 & 0 \\ -2 & 5 \end{smallmatrix}\bigr) = \bigl(\begin{smallmatrix} 4 + 2 & 0 - 5 \\ 2 - 6 & 0 + 15 \end{smallmatrix}\bigr) = \bigl(\begin{smallmatrix} 6 & -5 \\ -4 & 15 \end{smallmatrix}\bigr)$. Via $E F + E = \bigl(\begin{smallmatrix} 4 & -4 \\ -5 & 12 \end{smallmatrix}\bigr) + \bigl(\begin{smallmatrix} 2 & -1 \\ 1 & 3 \end{smallmatrix}\bigr) = \bigl(\begin{smallmatrix} 6 & -5 \\ -4 & 15 \end{smallmatrix}\bigr)$. Match.

    \item $E^{T} e_{1} = \bigl(\begin{smallmatrix} 2 \\ -1 \end{smallmatrix}\bigr)$ (first column of $E^{T}$, i.e.\ first row of $E$). $E^{T} e_{2} = \bigl(\begin{smallmatrix} 1 \\ 3 \end{smallmatrix}\bigr)$.

    \item $S^{2} = \bigl(\begin{smallmatrix} 0 & 1 \\ 1 & 0 \end{smallmatrix}\bigr)\bigl(\begin{smallmatrix} 0 & 1 \\ 1 & 0 \end{smallmatrix}\bigr) = \bigl(\begin{smallmatrix} 1 & 0 \\ 0 & 1 \end{smallmatrix}\bigr) = I_{2}$. ($S$ is an involution.)

    \item $S T = \bigl(\begin{smallmatrix} 0 & 1 \\ 1 & 0 \end{smallmatrix}\bigr)\bigl(\begin{smallmatrix} 1 & 2 \\ 3 & 4 \end{smallmatrix}\bigr) = \bigl(\begin{smallmatrix} 3 & 4 \\ 1 & 2 \end{smallmatrix}\bigr)$ (swaps the rows of $T$). $T S = \bigl(\begin{smallmatrix} 1 & 2 \\ 3 & 4 \end{smallmatrix}\bigr)\bigl(\begin{smallmatrix} 0 & 1 \\ 1 & 0 \end{smallmatrix}\bigr) = \bigl(\begin{smallmatrix} 2 & 1 \\ 4 & 3 \end{smallmatrix}\bigr)$ (swaps the columns of $T$).

    \item $R^{2} = \bigl(\begin{smallmatrix} 2 & 2 \\ 2 & 2 \end{smallmatrix}\bigr)\bigl(\begin{smallmatrix} 2 & 2 \\ 2 & 2 \end{smallmatrix}\bigr) = \bigl(\begin{smallmatrix} 8 & 8 \\ 8 & 8 \end{smallmatrix}\bigr) = 4 R$. By induction $R^{n} = 4^{n - 1} R$ for $n \ge 1$.

    \item $E - F = \bigl(\begin{smallmatrix} 1 & -1 \\ 3 & -1 \end{smallmatrix}\bigr)$, $E + F = \bigl(\begin{smallmatrix} 3 & -1 \\ -1 & 7 \end{smallmatrix}\bigr)$. $(E - F)(E + F) = \bigl(\begin{smallmatrix} 3 + 1 & -1 - 7 \\ 9 + 1 & -3 - 7 \end{smallmatrix}\bigr) = \bigl(\begin{smallmatrix} 4 & -8 \\ 10 & -10 \end{smallmatrix}\bigr)$. And $E^{2} - F^{2} = \bigl(\begin{smallmatrix} 3 & -5 \\ 5 & 8 \end{smallmatrix}\bigr) - \bigl(\begin{smallmatrix} 1 & 0 \\ -10 & 16 \end{smallmatrix}\bigr) = \bigl(\begin{smallmatrix} 2 & -5 \\ 15 & -8 \end{smallmatrix}\bigr) \neq (E - F)(E + F)$ in general (since $E F \neq F E$).

    \item $L^{2} = \bigl(\begin{smallmatrix} a & 0 \\ b & c \end{smallmatrix}\bigr)\bigl(\begin{smallmatrix} a & 0 \\ b & c \end{smallmatrix}\bigr) = \bigl(\begin{smallmatrix} a^{2} & 0 \\ a b + b c & c^{2} \end{smallmatrix}\bigr)$. $(1,2)$-entry is $0$: lower-triangular.

    \item If $A B = I$ for $2 \times 2$ matrices, take determinants: $\det A \cdot \det B = 1$, so both are nonzero, hence $A$ is invertible. Multiplying $A B = I$ on the left by $A^{-1}$: $B = A^{-1}$, and then $B A = A^{-1} A = I$. (For non-square or infinite-dimensional matrices this can fail.)

    \item $N^{T} N = \bigl(\begin{smallmatrix} 2 & 1 & -1 \\ 1 & 0 & 2 \end{smallmatrix}\bigr)\begin{pmatrix} 2 & 1 \\ 1 & 0 \\ -1 & 2 \end{pmatrix} = \bigl(\begin{smallmatrix} 4 + 1 + 1 & 2 + 0 - 2 \\ 2 + 0 - 2 & 1 + 0 + 4 \end{smallmatrix}\bigr) = \bigl(\begin{smallmatrix} 6 & 0 \\ 0 & 5 \end{smallmatrix}\bigr)$.
\end{solutions}
